\chapter{反抗我们自己}
《叛逆》真正让我们与自己为敌的时刻，大约出现在影片开头的四分之一处。
在那之前，影片呈现的是一个幸福世界。
任何还有一丝同理心的观众都不得不承认，这些角色早就配得上这样的幸福。
整部动画一路走来，我们跟随着这些少女，陪她们受苦，也在几乎无望的处境中盼她们终能获得幸福。

从某种意义上说，动画结尾确实给了她们这样的幸福，
但不可否认，那仍是一个苦乐参半的结局。
焰孤身一人，只有她还记得小圆。
小圆永远地离去了，甚至打从一开始就不曾出生。
其他魔法少女依旧战斗，依旧受苦，依旧会沉入绝望的深渊、化为魔女。
只是在她们化为魔女的一刻，小圆会以慈悲之手终结她们。

沙耶香仍旧为一个几乎没有在意过她的男孩而付出生命。
麻美和杏子仍作为魔法少女而战斗，这也意味着我们可以推断：
麻美的父母依旧已经去世。
杏子的家庭也毁于那场父亲杀人后自杀的惨剧。

动画的结尾是坦诚的，却不是幸福的。
它创造了一个新的、更好的世界，但那个世界远非完美。

而《叛逆》开篇呈现的，正是那个完美世界。
五位魔法少女全都活着，并肩作战。
她们之间的人际矛盾，被降格为沙耶香和杏子之间带着暧昧意味的打闹。
焰在反复时间旅行中承受的心理创伤似乎已经愈合：
她又变回了那个戴眼镜、扎双马尾、羞怯而不够自信，却更愉快、更温和的自己；
小圆也更加开朗、自信，更接近第十集开头焰初遇她时的模样。

至于梦魇，就其威胁而言，几乎显得滑稽可笑。
以仁美的梦魇为例，它们既不会对少女们造成物理上的伤害，也不会施加心理上的折磨，甚至还能被还原成字面意义上的“萌宠”。
更重要的是，梦魇被消灭后会释放大量净化灵魂宝石的光芒。
正如我先前指出的，这不仅允许少女们协同作战，甚至鼓励她们如此行事。
同时，这些光芒也提供了充足的能量，让她们不必担心灵魂宝石因污浊变黑而导致自己死亡。
于是，少女们得以协同施展强大的魔法，在华丽的视觉效果中迎战恰到好处的困难。
这既能给她们带来成就感，又使她们永远不必经历动画本篇中的那些恐怖。

这正是我们这些观众共同期待的。
诚然，大多数人都明白，原本的结局从美学上看或许更优秀。
可是数量庞大的同人创作已经说明了问题：
无论是因为那些创作者太缺乏经验而不懂，还是因为太过投入而不在乎，我们确实都希望这些女孩能拥有更美好的结局。
于是电影来了，给了我们想要的一切，直到焰开始察觉真相。

就像前文讨论过的《失乐园》（见第十九章）一样，影片诱使我们站到一个正在摧毁我们乐园的人那边。
焰知道这个幸福世界并不真实，因此我们也知道：
只要她继续追查，它就必然崩塌。
从她与杏子的第一次谈话开始，这个世界便越来越虚幻。
等到两人意识到自己被困在这座城市里时，周遭已经抽象成红色色块与白色线条，唯一还能辨认的物体只剩那辆公交车。
不久之后，焰又变回了我们在动画中熟悉的模样：
摘下眼镜，披着直发，阴郁而坚忍。
这个熟悉的形象令我们兴奋，可它同时也意味着，那个幸福世界正在以更快的速度瓦解。

随后，我们看到了在动画前三集就埋下伏笔的那场战斗：
麻美与焰终于兵戎相见。
这场战斗在视觉上令人惊叹，焰与麻美都将各自的能力和庞大的武器库发挥到极致。
它扣人心弦，充满戏剧张力，动画和配乐都堪称出色，可它又彻底不对劲。
作为一场动作大戏，它篇幅很长，却几乎不推进剧情，更没有起到塑造人物或深化主题的作用；
它当然令人血脉偾张，却明显是一段纯粹迎合观众的内容，而这正是动画本篇大多数时候刻意回避的事。
紧接着，焰朝自己的头开枪，麻美化作缎带，先前的讨好在一瞬间转为恐怖。

得偿所愿，带来的反而是失望，并最终通向恐怖。
影片高潮处对此表现得淋漓尽致：
焰与小圆重逢，一切却彻底失控。
最终诞生了一个所有少女都自由且活着、小圆也不必成为魔法少女的世界。
可那同时也是一个腐败的世界，由恶魔造物主般的焰统治，而小圆被她囚禁在其中。

我们买了票，签下契约，实现了愿望，而这些愿望却在我们身边化为灰烬。
欲望不可避免地通向痛苦。

为什么？
因为我们或许渴望幸福，但我们需要真相。
这不是说绝望比幸福更真实，而是说《魔圆》的真相在于熵，在于衰败是不可避免的。
而这部作品始终把物理意义上的熵与衰败，同抑郁和绝望的感受联系在一起。
若它干脆、单纯、毫不复杂地以幸福收场，又把我们想要的东西拱手献上，既不使之腐败，也不将其夺走，那反而是对作品自身的背弃。

因此，电影迫使我们拒绝自己对这个系列的愿望。
那些沉醉于作品黑暗与棘手难度的观众，必须忍受“杀必死\footnote{
    译注：网络流行语，原意是指代二次元作品中的“福利桥段”或“擦边球”内容，此处指作品为迎合粉丝而设计的画面。
}”的奉承，从而被迫否认自己的粉丝身份。
而那些愿意拥抱“杀必死”的观众，则必须面对它最终通向何处。
双方都必须直面影片最后一段极其暧昧的走向：
焰创造了一个比动画本篇更黑暗，也更明亮的世界，而这个世界又比影片前半部分的梦境世界更加连贯自洽。

由此，作品把观众放到了魔法少女的位置上。
一味要求这个系列实现我们的愿望，最终会使它走向悲剧。
愿望变成诅咒，让我们一步步走向深渊。
直到影片高潮，我们竟在期待身为魔女的焰能够撕裂世界。
而当她真的这么做了之后，我们又必须面对这个愿望所创造出的现实。

电影挑动我们与自身为敌，并让我们看见自己对这个系列的愿望是如何背叛自己的。
借此，它最后一次试图把同理心推到我们面前。
就像动画最初几集那样，它先以视觉盛宴和“杀必死”吸引我们。
等陷阱布好、诱饵就位，它便迫使我们与角色感同身受。
更进一步说，它不是让我们同情角色，而是让我们以类似角色的方式去共情。
焰把小圆剥离的那一刻，那种涌上心头的困惑、疏离和强烈的违和感，正是成为魔法少女是什么滋味的一小片回声。

我之前提到过，这部作品关乎熵与衰败，关乎抑郁与绝望。
事实也确实如此。
但我们也很容易忘记，它还关乎别的东西。
而通过让我们与自身为敌，影片将那些东西重新带回到我们眼前。

它不只是关于佛教、德国文学或魔法少女题材。
它不只是关于熵与苦难，不只是关于主题的复杂性，也不只是关于那个隐含的、格式塔式作者可能具有的心理问题。
它甚至不只是关于角色，不只是那些由动画师创造、由声优赋声的光影与色块。
它也关乎我们：
关于我们的恐惧与希望，关于我们的愿望与绝望，关于我们共情他人的能力。

归根结底，从一开始，《魔法少女小圆》就是一个关于人的故事。
